% Abhakseite „Das kann ich" – Zone nur im Gesamt (\mitzone)
\begin{abhakseite}
\ifdefined\mitzone
\abhakgruppe{Kennst du schon}
\abhak{1}{Ich kann Zahlen quadrieren, auch negative Zahlen und Dezimalzahlen.}
\abhak{2}{Ich kann mit Klammern und negativen Zahlen rechnen: Klammer zuerst, Punkt vor Strich.}
\abhak{3}{Ich kann eine lineare Gleichung lösen.}
\abhak{4}{Ich kann durch Einsetzen prüfen, ob eine Zahl eine Lösung ist.}
\abhak{5}{Ich kann Quadratwurzeln ziehen.}
\abhak{6}{Ich kann den Scheitel aus der Scheitelpunktform ablesen.}
\abhak{7}{Ich kann Klammern ausmultiplizieren, auch mit den binomischen Formeln.}
\abhak{8}{Ich kann ausklammern.}
\abhak{9}{Ich kann Terme ordnen und zusammenfassen.}
\abhak{10}{Ich finde den Fehler beim Auflösen einer quadrierten Klammer.}
\abhak{11}{Ich kann eine quadrierte Klammer auflösen und zusammenfassen.}
\fi
\abhakgruppe{Einheit 1 · Wurzelziehen und Lösbarkeit}
\abhak{12}{Ich erkenne, ob eine Gleichung quadratisch ist.}
\abhak{13}{Ich erkenne, ob rechts eine positive Zahl, null oder eine negative Zahl steht.}
\abhakgruppe{Einheit 1 · Gleichungen der Form $x^2 = c$}
\abhak{14}{Ich kann eine Gleichung $x^2 = c$ durch Wurzelziehen lösen.}
\abhak{15}{Ich kann $x^2$ erst allein stellen und dann die Wurzel ziehen.}
\abhakgruppe{Einheit 1 · Gleichungen mit quadrierter Klammer}
\abhak{16}{Ich kann eine Gleichung mit quadrierter Klammer rückwärts lösen.}
\abhakgruppe{Einheit 1 · Lösungen prüfen und zählen}
\abhak{17}{Ich kann prüfen, ob eine Zahl eine quadratische Gleichung löst.}
\abhak{18}{Ich kann am Graphen ablesen, wie viele Lösungen eine Gleichung hat.}
\abhak{19}{Ich kann Gleichungen mit keiner, einer oder zwei Lösungen angeben.}
\abhak{20}{Ich kann Aussagen über die Lösungen einer Gleichung beurteilen.}
\abhak{21}{Ich finde den Fehler beim Wurzelziehen.}
\abhak{22}{Ich kann begründen, wie viele Lösungen eine Gleichung hat.}
\abhak{23}{Ich kann mit Wurzelziehen eine Länge berechnen.}
\abhakgruppe{Einheit 2 · Normalform und p-q-Formel}
\abhak{24}{Ich erkenne an der Form einer Gleichung, welcher Lösungsweg passt.}
\abhak{25}{Ich erkenne p und q in einer Gleichung in Normalform.}
\abhak{26}{Ich erkenne, ob vor $x^2$ eine Eins steht.}
\abhakgruppe{Einheit 2 · Lösen mit der p-q-Formel}
\abhak{27}{Ich kann eine Gleichung in Normalform mit der p-q-Formel lösen. (je nach Buch erst in Klasse 10)}
\abhak{28}{Ich kann eine Gleichung in Normalform mit der p-q-Formel lösen – weiter.}
\abhakgruppe{Einheit 2 · Gleichungen ordnen und durch die Zahl vor $x^2$ teilen}
\abhak{29}{Ich kann eine Gleichung erst ordnen und dann mit der p-q-Formel lösen.}
\abhak{30}{Ich kann eine Gleichung mit einer Zahl vor $x^2$ lösen.}
\abhak{31}{Ich kann eine Gleichung mit Klammern erst auflösen, ordnen und dann lösen.}
\abhakgruppe{Einheit 2 · Zahl der Lösungen}
\abhak{32}{Ich kann eine Gleichung mit genau einer Lösung lösen.}
\abhak{33}{Ich kann erkennen, dass eine Gleichung keine Lösung hat.}
\abhak{34}{Ich kann am Wert unter der Wurzel entscheiden, wie viele Lösungen es gibt.}
\abhakgruppe{Einheit 2 · Probe und Lösungsweg}
\abhak{35}{Ich kann mit einer Probe prüfen, ob eine Lösung stimmt.}
\abhak{36}{Ich kann den kürzesten Lösungsweg wählen.}
\abhak{37}{Ich finde den Fehler beim Lösen mit der p-q-Formel.}
\abhak{38}{Ich kann begründen, wann die p-q-Formel funktioniert.}
\abhak{39}{Ich kann berechnen, wo sich eine Gerade und eine Parabel schneiden.}
\abhakgruppe{Einheit 3 · Sachaufgaben}
\abhak{40}{Ich erkenne, welche Lösung zur Frage passt.}
\abhakgruppe{Einheit 3 · Zahlenrätsel}
\abhak{41}{Ich kann ein Zahlenrätsel mit dem Quadrat einer Zahl lösen.}
\abhak{42}{Ich kann zu einem Zahlenrätsel selbst die Gleichung aufstellen.}
\abhakgruppe{Einheit 3 · Flächenaufgaben}
\abhak{43}{Ich kann eine Rechteckaufgabe mit einer quadratischen Gleichung lösen.}
\abhak{44}{Ich kann eine Flächenaufgabe mit Klammern lösen.}
\abhak{45}{Ich finde den Fehler in einer Sachaufgabe.}
\abhak{46}{Ich kann begründen, welche Lösungen im Sachzusammenhang gelten.}
\abhakgruppe{Ausblick Einheit 4 · Satz vom Nullprodukt}
\abhak{47}{Ich erkenne, ob rechts eine Null steht.}
\abhakgruppe{Ausblick Einheit 4 · Produkt gleich null}
\abhak{48}{Ich kann eine Gleichung in Produktform mit dem Satz vom Nullprodukt lösen.}
\abhak{49}{Ich kann prüfen, welche Zahl eine Gleichung in Produktform erfüllt.}
\abhakgruppe{Ausblick Einheit 4 · Ausklammern}
\abhak{50}{Ich kann eine Gleichung ohne Zahlglied durch Ausklammern lösen.}
\abhakgruppe{Ausblick Einheit 4 · Produkt ungleich null}
\abhak{51}{Ich kann eine Produktgleichung lösen, bei der rechts nicht null steht.}
\abhak{52}{Ich kann eine Gleichung lösen, indem ich die linke Seite als Produkt schreibe.}
\abhak{53}{Ich finde den Fehler beim Satz vom Nullprodukt.}
\abhak{54}{Ich kann begründen, wann der Satz vom Nullprodukt gilt.}
\abhak{55}{Ich kann mit dem Satz vom Nullprodukt eine Sachfrage beantworten.}
\end{abhakseite}
